\documentclass[answers,12pt,addpoints]{exam} 
\usepackage{import}
\usepackage{xcolor}
\usepackage{amsmath}
\usepackage{amsfonts}
\usepackage{amssymb}

% Define blue text for problematic areas - FIXED for multi-paragraph content
\newcommand{\issue}[1]{{\color{blue}#1}}
\newcommand{\warning}[1]{{\color{red}\textbf{WARNING:} #1}}

% Header
\newcommand{\name}{Pranav Tikkawar}
\newcommand{\course}{Spatial Statistics}
\newcommand{\assignment}{Research}
\author{\name}
\title{\course \\ - \assignment}

\begin{document}
\maketitle

\section{A General Periodic Fourier--State-Space Spatio--Temporal Model}

\subsection{Model setup}

We consider a real-valued Gaussian spatio--temporal process
\[
  \{X_t(x) : x \in [0,2\pi],\; t \in \mathbb{Z}\},
\]
defined on a one-dimensional spatial domain with periodic boundary conditions (a circle) and discrete time. For each time $t$, the spatial field $X_t(\cdot)$ is represented by a truncated Fourier series
\begin{equation}
  X_t(x)
  \;=\;
  a_0(t)
  \;+\;
  \sum_{j=1}^{J}
    \Bigl[
      a_j(t)\cos(jx) + b_j(t)\sin(jx)
    \Bigr],
  \qquad x \in [0,2\pi],
  \label{eq:Xt-fourier}
\end{equation}
where $J$ is a truncation level and
\[
  a_j(t),\, b_j(t), \quad j=0,\dots,J,\; t\in\mathbb{Z},
\]
are jointly Gaussian random coefficients.

For $j \ge 1$, the pair $(a_j(t), b_j(t))$ corresponds to the $j$th spatial frequency (wavenumber). Low $j$ terms encode large-scale, slowly varying spatial structure, while high $j$ terms encode fine-scale spatial detail.

\subsection{Spatial margin: Matérn-like variances}

The spatial covariance at each fixed time is controlled by a Matérn-like spectral density in the Fourier domain. For $j = 0,\dots,J$ we specify
\begin{equation}
  \sigma_j^2
  \;\propto\;
  \bigl(\kappa^2 + j^2\bigr)^{-(\nu + 1/2)},
  \quad
  \kappa = \frac{\sqrt{2\nu}}{\ell},
  \label{eq:matern-spectrum}
\end{equation}
where
\begin{itemize}
  \item $\nu > 0$ is the Matérn smoothness parameter,
  \item $\ell > 0$ is a spatial range parameter,
  \item $\kappa = \sqrt{2\nu}/\ell$ is an inverse scale parameter.
\end{itemize}

{\color{blue}
NOTE: This spectral form is appropriate for a \emph{non-periodic} domain $\mathbb{R}$. On a periodic domain $[0,2\pi]$, the correspondence between this discrete spectrum and the continuous Matérn model requires careful treatment. 

\textbf{Question:} Should we (a) sample the continuous Matérn spectrum at integer frequencies $j$, or (b) explicitly periodize the Matérn covariance function and compute its Fourier coefficients?
}

For simulation we normalize these variances so that the total marginal variance of $X_t(x)$ is $\sigma^2_{\text{tot}}$. Defining
\[
  \tilde{\sigma}_j^2
  \;=\;
  \bigl(\kappa^2 + j^2\bigr)^{-(\nu + 1/2)},
\]
we set
\begin{equation}
  \sigma_j^2
  \;=\;
  \sigma^2_{\text{tot}}
  \cdot
  \frac{\tilde{\sigma}_j^2}{\sum_{k=0}^{J} \tilde{\sigma}_k^2},
  \qquad j = 0,\dots,J.
  \label{eq:sigma-normalization}
\end{equation}

\subsection{Temporal dynamics: block-diagonal AR(1) with rotation and correlated innovations}

We collect the Fourier coefficients into a state vector
\begin{equation}
  \mathbf{Z}_t
  \;=\;
  \bigl(a_0(t), a_1(t),\dots,a_J(t), b_1(t),\dots,b_J(t)\bigr)^\top
  \;\in\; \mathbb{R}^{2J+1}.
  \label{eq:state-vector}
\end{equation}

The temporal evolution is governed by a \emph{block-diagonal} linear Gaussian state-space model. For the constant mode $j=0$, we have a scalar AR(1):
\begin{equation}
  a_0(t+1) = \rho_0\, a_0(t) + \eta_{0,t},
  \qquad \eta_{0,t} \sim \mathcal{N}\bigl(0, \sigma_0^2(1 - \rho_0^2)\bigr),
  \label{eq:ar1-j0}
\end{equation}
with $|\rho_0| < 1$ ensuring stationarity.

For each mode $j \ge 1$, the cosine and sine coefficients evolve jointly according to a $2 \times 2$ block AR(1):
\begin{equation}
  \begin{pmatrix} a_j(t+1) \\ b_j(t+1) \end{pmatrix}
  =
  M_j
  \begin{pmatrix} a_j(t) \\ b_j(t) \end{pmatrix}
  +
  \begin{pmatrix} \eta_{a,j,t} \\ \eta_{b,j,t} \end{pmatrix},
  \label{eq:ar1-block}
\end{equation}
where the $2 \times 2$ transition matrix is
\begin{equation}
  M_j
  =
  \rho_j
  \begin{pmatrix}
    \cos\theta_j & -\sin\theta_j \\
    \sin\theta_j & \cos\theta_j
  \end{pmatrix},
  \label{eq:Mj-rotation}
\end{equation}
with $\rho_j \in (0,1)$ (amplitude persistence) and $\theta_j \in \mathbb{R}$ (rotation angle).

The $2 \times 2$ innovation covariance for mode $j$ is
\begin{equation}
  \begin{pmatrix} \eta_{a,j,t} \\ \eta_{b,j,t} \end{pmatrix}
  \sim
  \mathcal{N}\Bigl(\mathbf{0},\; \Sigma_{\varepsilon,j}\Bigr),
  \qquad
  {\color{blue}\Sigma_{\varepsilon,j}
  =
  \sigma_j^2(1 - \rho_j^2)
  \begin{pmatrix} 1 & \phi_j \\ \phi_j & 1 \end{pmatrix}},
  \label{eq:innovation-cov}
\end{equation}
where $\phi_j \in (-1,1)$ is the \emph{within-frequency correlation parameter} between the cosine and sine innovations.

{\color{red}
\textbf{WARNING — EMPIRICAL VALIDATION ISSUE:} The Lyapunov equation $\Sigma_j = M_j \Sigma_j M_j^\top + \Sigma_{\varepsilon,j}$ with target stationary covariance $\Sigma_j = \sigma_j^2 \begin{pmatrix} 1 & \phi_j \\ \phi_j & 1 \end{pmatrix}$ does yield \eqref{eq:innovation-cov} as shown in Section \ref{sec:lyapunov-fix}. However, empirical validation showed:
\begin{itemize}
  \item Empirical $\operatorname{Cor}(a_1, b_1) = 0.208$, theoretical $\phi_1 = 0.696$ (70\% error)
  \item Phase drift empirically $\approx 0$ even when $\theta_j \neq 0$ was set
  \item ACF of amplitude decays faster than $\rho_j^k$ (up to 57\% error for high $j$)
\end{itemize}
This suggests either (a) an implementation error in the R code, or (b) a misunderstanding of what quantity should match theory when rotation is present.
}

The innovations are independent across different frequencies $j$ and across time $t$.

\subsubsection*{Physical interpretation}

\paragraph{Rotation parameter $\theta_j$:}
When $\theta_j = 0$, the transition matrix $M_j = \rho_j I_2$, so cosine and sine evolve independently. When $\theta_j \neq 0$, we have temporal mixing:
\[
  a_j(t+1) = \rho_j[\cos\theta_j \cdot a_j(t) - \sin\theta_j \cdot b_j(t)] + \eta_{a,j,t},
\]
inducing \emph{phase rotation}: spatial patterns propagate around the circle over time. The eigenvalues of $M_j$ are $\rho_j e^{\pm i\theta_j}$, so stationarity holds for $|\rho_j| < 1$ regardless of $\theta_j$.

{\color{blue}
\textbf{Question:} Does phase rotation $\theta_j$ have a physical interpretation in terms of wave propagation speed on the circle? If $\theta_j = c \cdot j$ for constant $c$, does this correspond to a traveling wave with speed $c$?
}

\paragraph{Innovation correlation parameter $\phi_j$:}
When $\phi_j = 0$, the innovations $\eta_{a,j,t}$ and $\eta_{b,j,t}$ are independent. When $\phi_j \neq 0$, the innovations are correlated: $\operatorname{Cov}(\eta_{a,j,t}, \eta_{b,j,t}) = \sigma_j^2(1-\rho_j^2)\phi_j$. {\color{blue}At stationarity, this targets $\operatorname{Cov}(a_j(t), b_j(t)) = \sigma_j^2 \phi_j$ via the Lyapunov equation (Section \ref{sec:lyapunov-fix}).}

Typical parameterization: $\phi_j = \phi_0 / (1 + j)$ or $\phi_j = \phi_0 e^{-\gamma j}$.

\paragraph{Combined effect:}
$\theta_j \neq 0$ creates \emph{temporal cross-dependence}; $\phi_j \neq 0$ creates \emph{contemporaneous correlation}. Together they allow rich within-frequency dynamics while maintaining the block-diagonal structure.

{\color{blue}
\textbf{Key empirical finding:} Simulated $\operatorname{Cor}(a_j, b_j)$ showed 70\% error for $j=1$ and wild scatter for other modes, suggesting either incorrect initial condition draws or a numerical issue in the multivariate normal generation.
}

\subsection{Stationary covariance of $(a_j(t), b_j(t))$ via Lyapunov equation}
\label{sec:lyapunov-fix}

At stationarity, the $2\times2$ covariance matrix
\[
  \Sigma_{j}
  =
  \operatorname{Cov}\begin{pmatrix}a_j(t)\\b_j(t)\end{pmatrix}
\]
satisfies the discrete-time Lyapunov equation
\begin{equation}
  \Sigma_j = M_j \Sigma_j M_j^\top + \Sigma_{\varepsilon,j}.
  \label{eq:lyapunov-2x2}
\end{equation}

{\color{blue}
\textbf{Corrected derivation:} For the rotation-scaling matrix \eqref{eq:Mj-rotation}, we have
\[
  M_j M_j^\top = \rho_j^2 I_2,
\]
so \eqref{eq:lyapunov-2x2} becomes
\[
  \Sigma_j = \rho_j^2 \Sigma_j + \Sigma_{\varepsilon,j}
  \quad \Longrightarrow \quad
  \Sigma_{\varepsilon,j} = (1 - \rho_j^2) \Sigma_j.
\]
If we target the stationary covariance
\[
  \Sigma_j = \sigma_j^2 \begin{pmatrix} 1 & \phi_j \\ \phi_j & 1 \end{pmatrix},
\]
then the innovation covariance must be
\[
  \Sigma_{\varepsilon,j} = \sigma_j^2(1 - \rho_j^2) \begin{pmatrix} 1 & \phi_j \\ \phi_j & 1 \end{pmatrix}.
\]
This formula is \emph{independent of} $\theta_j$ because $M_j M_j^\top = \rho_j^2 I_2$ for any rotation angle. Therefore equation \eqref{eq:innovation-cov} is theoretically correct, and the empirical mismatch must arise from implementation or interpretation errors.
}

{\color{blue}
\textbf{Advisor questions:}
\begin{enumerate}
  \item The empirical ACF of amplitude $R_j(t) = \sqrt{a_j(t)^2 + b_j(t)^2}$ decays much faster than $\rho_j^k$ (up to 57\% error for high $j$). Should the theoretical ACF be computed for the \emph{amplitude} instead of the coefficients?
  
  \item The empirical correlation $\operatorname{Cor}(a_j(t), b_j(t))$ showed wild scatter instead of matching $\phi_j \approx 0.65$--0.70. Is this a finite-sample effect, or is there a bias in how the stationary initial condition is drawn?
  
  \item The spatial covariance at large lags showed an ``upturn'' (U-shape) instead of monotonic decay. Is this the ``ridge'' artifact from nonseparable models warned about in Stein (2005)? Does this indicate we need different temporal structure?
\end{enumerate}
}

Therefore, at stationarity:
\begin{align}
  \operatorname{Var}(a_j(t)) &= \sigma_j^2, \\
  \operatorname{Var}(b_j(t)) &= \sigma_j^2, \\
  \operatorname{Cov}(a_j(t), b_j(t)) &= \sigma_j^2 \phi_j.
\end{align}

The stationary covariance depends on $\phi_j$ but \emph{not} on $\theta_j$. The rotation angle affects temporal evolution but not marginal variances or contemporaneous correlation.

\subsection{Nonseparability via frequency-dependent persistence}

The key nonseparable feature arises from allowing $\rho_j$ to depend on spatial frequency $j$. We propose an \emph{algebraic (power-law) decay} parameterization
\begin{equation}
  \rho_j
  =
  \frac{1}{(1 + \lambda j^\alpha)^\beta},
  \qquad j = 1,\dots,J,
  \label{eq:rhoj-algebraic}
\end{equation}
with $\rho_0 \in (0,1)$ set separately, and parameters $\lambda > 0$, $\alpha > 0$, $\beta > 0$.


Under this structure, low-frequency modes have $\rho_j \approx 1$ (slow temporal decay), while high-frequency modes have smaller $\rho_j$ (fast decay). The effective spatial correlation structure changes with time lag: at short lags both low and high frequencies contribute, but at long lags only low frequencies remain. This space--time covariance cannot be written as $C_{\text{space}}(h) C_{\text{time}}(k)$ and is therefore nonseparable.

\subsection{Space--time covariance}

Under the block-diagonal structure, the space--time covariance
\[
  C(h,k) = \mathbb{E}\bigl[ X_t(x) X_{t+k}(x+h) \bigr]
\]
can be computed when different modes are uncorrelated across frequencies. For $\theta_j = 0$ (no rotation), we have
\[
  \mathbb{E}[a_j(t)a_j(t+k)] = \mathbb{E}[b_j(t)b_j(t+k)] = \sigma_j^2 \rho_j^{|k|},
\]
yielding
\begin{equation}
  C(h,k)
  =
  \sum_{j=0}^{J}
  \sigma_j^2 \rho_j^{|k|} \cos(jh).
  \label{eq:cov-fourier-ar1}
\end{equation}

{\color{blue}
\textbf{Derivation assumes $\theta_j = 0$:} When $M_j$ includes rotation, the lag-$k$ covariance involves $M_j^k$:
\[
  M_j^k = \rho_j^k \begin{pmatrix} \cos(k\theta_j) & -\sin(k\theta_j) \\ \sin(k\theta_j) & \cos(k\theta_j) \end{pmatrix},
\]
so
\[
  \mathbb{E}[a_j(t) a_j(t+k)] = \sigma_j^2 \rho_j^k \cos(k\theta_j) + \text{cross-terms involving } \phi_j.
\]
The spatial covariance \eqref{eq:cov-fourier-ar1} is therefore correct \emph{only when} $\theta_j = 0$ or when cross-terms cancel in the Fourier reconstruction. \textbf{please verify when rotation is present.}
}

\subsection{Parameter overview}

The full model is specified by the following parameters, organized by their role in the model structure.

\subsubsection*{Fourier truncation}

\begin{itemize}
  \item $J \in \mathbb{Z}_+$: highest Fourier mode (truncation level). Controls the finest spatial scale represented. Larger $J$ allows finer spatial detail; in practice $J$ is chosen so that the total variance in modes $j > J$ is negligible.
\end{itemize}

\subsubsection*{Spatial margin parameters (Matérn-like spectrum)}

\begin{itemize}
  \item $\nu > 0$: Matérn smoothness parameter. Larger $\nu$ gives smoother spatial fields. Sample paths are $\lceil \nu \rceil - 1$ times mean-square differentiable.
  \item $\ell > 0$: spatial range parameter. Larger $\ell$ increases the spatial correlation length by shifting spectral mass to lower frequencies.
  \item $\sigma^2_{\text{tot}} > 0$: total marginal variance of the field $X_t(x)$. Overall scale parameter.
  \item $\kappa = \sqrt{2\nu}/\ell$: (derived) inverse range parameter used in the spectral density \eqref{eq:matern-spectrum}.
\end{itemize}

The mode variances $\sigma_j^2$ for $j=0,\dots,J$ are derived from $(\nu, \ell, \sigma^2_{\text{tot}})$ via \eqref{eq:matern-spectrum} and \eqref{eq:sigma-normalization}.

\subsubsection*{Temporal persistence parameters}

\begin{itemize}
  \item $\rho_0 \in (0,1)$: AR(1) persistence for the constant mode $a_0(t)$. Set separately to avoid nonstationarity (the formula \eqref{eq:rhoj-algebraic} would give $\rho_0 = 1$ for $j=0$).
  \item $\lambda > 0$: controls overall temporal decay strength across frequencies in \eqref{eq:rhoj-algebraic}.
  \item $\alpha > 0$: controls how strongly persistence changes with frequency $j$ in \eqref{eq:rhoj-algebraic}.
  \item $\beta > 0$: decay exponent in \eqref{eq:rhoj-algebraic}. Smaller $\beta$ gives slower (longer-memory) decay.
\end{itemize}

The mode-specific persistence parameters $\rho_j$ for $j=1,\dots,J$ are derived from $(\lambda, \alpha, \beta)$ via \eqref{eq:rhoj-algebraic}.

\subsubsection*{Phase rotation parameters}

\begin{itemize}
  \item $\theta_j \in \mathbb{R}$ for $j=1,\dots,J$: rotation angle for mode $j$ in the $2\times2$ transition matrix $M_j$ (equation \eqref{eq:Mj-rotation}).
  \item When $\theta_j = 0$ for all $j$, the model reduces to independent cosine/sine evolution (no temporal mixing).
  \item Typical parameterization: $\theta_j = \theta_0 / j$ or $\theta_j = \theta_0 / (1 + j)$ for a scalar hyperparameter $\theta_0 \in \mathbb{R}$, so that lower frequencies rotate more than higher frequencies.
\end{itemize}

{\color{blue}
\textbf{Empirical finding:} Phase drift validation showed mean $\Delta \text{phase} \approx 0$ even when $\theta_j$ was set to nonzero values. This suggests either (a) rotation is not being applied correctly in simulation, or (b) the phase of amplitude $R_j(t) = \sqrt{a_j^2 + b_j^2}$ does not drift at rate $\theta_j$ as expected. \textbf{What is the theoretical expected phase drift rate when rotation is present?}
}

\subsubsection*{Innovation correlation parameters}

\begin{itemize}
  \item $\phi_j \in (-1,1)$ for $j=1,\dots,J$: within-frequency correlation between innovations $\eta_{a,j,t}$ and $\eta_{b,j,t}$ (equation \eqref{eq:innovation-cov}).
  \item When $\phi_j = 0$ for all $j$, the innovations are independent (cosine and sine uncorrelated at the same time).
  \item When $\phi_j \neq 0$, the model induces contemporaneous correlation $\operatorname{Cov}(a_j(t), b_j(t)) = \sigma_j^2 \phi_j$ at stationarity.
  \item Typical parameterization: $\phi_j = \phi_0 / (1 + j)$ or constant $\phi_j = \phi_0$ for all $j$, where $\phi_0 \in (-1,1)$ is a scalar hyperparameter.
\end{itemize}

{\color{blue}
\textbf{Empirical finding:} Simulated $\operatorname{Cor}(a_j, b_j)$ showed 70\% error for $j=1$ (empirical 0.208 vs. theoretical 0.696) and high scatter for other modes. This suggests either (a) the stationary initial condition is drawn incorrectly, or (b) there is a numerical issue in the multivariate normal generation. The Lyapunov theory is correct (Section \ref{sec:lyapunov-fix}), so the issue is implementation-related.
}

\subsubsection*{Time horizon}

\begin{itemize}
  \item $T \in \mathbb{Z}_+$: number of time points to simulate or observe. Not a model parameter but a data specification.
\end{itemize}

\subsubsection*{Parameter count summary}

For a model with Fourier truncation $J$ and the simplest scalar hyperparameter choices, the total number of free parameters is:

\begin{itemize}
  \item \textbf{Minimal model} (independent modes, no rotation, no innovation correlation): $J=50$ modes
    \[
      \text{Parameters: } \nu, \ell, \sigma^2_{\text{tot}}, \rho_0, \lambda, \alpha, \beta
      \quad \Rightarrow \quad 7 \text{ parameters total.}
    \]
  
  \item \textbf{With rotation} (add $\theta_0$ for $\theta_j = \theta_0/j$):
    \[
      \text{Parameters: } 7 + 1 = 8.
    \]
  
  \item \textbf{With rotation and innovation correlation} (add $\theta_0$ and $\phi_0$ for $\phi_j = \phi_0/(1+j)$):
    \[
      \text{Parameters: } 7 + 1 + 1 = 9.
    \]
  
  \item \textbf{Fully flexible} (separate $\theta_j$ and $\phi_j$ for each mode $j=1,\dots,J$):
    \[
      \text{Parameters: } 7 + J + J = 7 + 2J.
    \]
    For $J=50$, this is 107 parameters—likely overparameterized for most applications.
\end{itemize}

{\color{blue}
\textbf{Question:} For inference (e.g., Whittle likelihood in Fourier space), which parameterization strikes the best balance between flexibility and identifiability? Should we use mode-specific $\{\rho_j, \theta_j, \phi_j\}$ as free parameters, or constrain them via smooth functions of $j$ with scalar hyperparameters?
}

By setting $\theta_j = 0$ and $\phi_j = 0$ for all $j$, the model reduces to the simple independent-mode AR(1) with uncorrelated cosine and sine components (7 parameters). Nonzero $\theta_j$ and $\phi_j$ introduce richer within-frequency dynamics: temporal mixing and phase rotation ($\theta_j$), and instantaneous correlation ($\phi_j$).

\section*{Open Questions}

\begin{enumerate}
  \item \textbf{Periodic Matérn spectrum:} Should we sample the continuous Matérn spectral density at integer frequencies $j$, or explicitly periodize the Matérn covariance and compute its Fourier series?
  
  \item \textbf{Rotation and ACF:} When $M_j$ includes rotation ($\theta_j \neq 0$), what is the theoretical ACF of amplitude $R_j(t) = \sqrt{a_j^2 + b_j^2}$? Does it match $\rho_j^k$, or does rotation change the decay rate?
  
  \item \textbf{Innovation correlation identifiability:} The empirical $\operatorname{Cor}(a_j, b_j)$ showed high variance. Is $\phi_j$ identifiable from data, or does it trade off with $\theta_j$?
  
  \item \textbf{Phase drift interpretation:} If $\theta_j = c \cdot j$, does this correspond to a traveling wave with speed $c$? Physical meaning of frequency-dependent rotation?
  
  \item \textbf{Cross-frequency dependence:} Current model has block-diagonal dynamics. Should we add off-block-diagonal elements in $M$ or $\Sigma_\varepsilon$ for cross-frequency coupling? What structure preserves computational efficiency?
  
  \item \textbf{Parameter parsimony for inference:} For Whittle likelihood estimation, should we use mode-specific $\{\rho_j, \theta_j, \phi_j\}$ as free parameters, or constrain via smooth functions with scalar hyperparameters (e.g., $\theta_j = \theta_0/j$)? What is the trade-off between flexibility and identifiability?
\end{enumerate}

\section{Next steps}
Write the Covariance in closed form as then it nice to simulate with fft due to circulant structure. 

If we truncate to $J$ it depends on how dense the items are

Want $Z(x) = \sum a_j(t) \cos(jx) + b_j(t) \sin(jx)$ to be a Gaussian process with covariance $C(h,k)$.

Consider Aliasing 
$Z(k/n) = \sum a_j(t) \cos(2\pi jk/n) + b_j(t) \sin(2\pi jk/n)$

Now simulat $A_j' = \sum_p A_{a_j+pn}$ and $B_j' = \sum_p B_{b_j+pn}$

and $\text{Var}(Z(k/n)) = \sum_{p = -\infty}^{\infty} \sigma^2_{j+pn} = \sum_{p=-q}^{q} \text{exact} + \sum_{p=q+1}^{\infty} + \sum_{p=-\infty} \text{tail by approximation}$


\textbf{Ask} about the proof for the other of psd is $\int_d$

Can use FFT for periodic statiorya. Paper: Chan \& Wood, Deitrich \& Neuwsam

maybe doing a smooth extenstion.

EQ 9 in stien paper JASA 2005

$Z(x,t+1) = \int v(x-y_) dy + \epsilon$




\end{document}